% RQ1 Methodology: Participation Dynamics Experiments

\subsection{Experimental Design}

We conduct two complementary experiments to investigate participation dynamics:

\begin{table}[htbp]
\centering
\caption{RQ1 Experiment configurations}
\label{tab:rq1-experiments}
\begin{tabular}{@{}llll@{}}
\toprule
Experiment & Focus & Swept Parameters & Runs \\
\midrule
01-Calibration & Participation targets & Target rate, incentive strength & 270 \\
03-Quorum & Quorum sensitivity & Quorum percentage (0.01--0.10) & 150 \\
\bottomrule
\end{tabular}
\end{table}

\subsection{Experiment 01: Participation Calibration}

\subsubsection{Objective}

Determine how participation target rates affect realized turnout and long-term voter retention.

\subsubsection{Parameter Grid}

We sweep a $3 \times 3$ grid:

\begin{itemize}
    \item \textbf{Participation target rate}: \{0.05, 0.15, 0.30\} (5\%, 15\%, 30\%)
    \item \textbf{Incentive multiplier}: \{0.5, 1.0, 2.0\} (low, medium, high)
\end{itemize}

Each configuration runs 30 times with different random seeds.

\subsubsection{Hypotheses}

\begin{itemize}
    \item \textbf{H1.1}: Higher participation targets lead to lower quorum achievement rates
    \item \textbf{H1.2}: Stronger incentives partially compensate for aggressive targets
    \item \textbf{H1.3}: Voter retention decreases when targets exceed natural participation capacity
\end{itemize}

\subsection{Experiment 03: Quorum Sensitivity}

\subsubsection{Objective}

Map the relationship between quorum thresholds and governance outcomes.

\subsubsection{Parameter Sweep}

Quorum percentage varied: 1\%, 2\%, 3\%, 4\%, 5\%, 6\%, 7\%, 8\%, 9\%, 10\%

Each threshold runs 15 times with different seeds.

\subsubsection{Hypotheses}

\begin{itemize}
    \item \textbf{H3.1}: Quorum-pass rate relationship is non-linear (S-curve)
    \item \textbf{H3.2}: Optimal quorum exists balancing legitimacy and operability
    \item \textbf{H3.3}: High quorums cause fatigue accumulation as agents strain to participate
\end{itemize}

\subsection{Baseline Configuration}

Both experiments use a common baseline:

\begin{itemize}
    \item Members: 200 agents
    \item Token distribution: Power-law ($\alpha = 1.5$)
    \item Agent mix: 5\% whales, 10\% delegates, 40\% active, 45\% passive
    \item Simulation length: 2,000 steps (days)
    \item Proposal frequency: 0.5/day average
    \item Voting period: 7 days
    \item Fatigue parameters: $\alpha=0.95$, $\beta=0.05$, $\gamma=0.01$
\end{itemize}

\subsection{Metrics}

Primary metrics for RQ1:

\begin{description}
    \item[Average Turnout] Mean participation rate across all proposals
    \item[Quorum Achievement Rate] Fraction of proposals reaching quorum
    \item[Voter Retention] Fraction of initially-active voters still participating at simulation end
    \item[Turnout Trend] Slope of turnout over time (decay rate)
    \item[Fatigue Accumulation] Mean fatigue level at simulation end
\end{description}

\subsection{Statistical Analysis}

\subsubsection{Summary Statistics}

For each parameter configuration:
\begin{itemize}
    \item Mean, median, standard deviation
    \item 95\% confidence intervals via bootstrap
    \item Interquartile range
\end{itemize}

\subsubsection{Hypothesis Testing}

\begin{itemize}
    \item \textbf{ANOVA}: Compare means across participation target levels
    \item \textbf{Regression}: Model turnout as function of quorum threshold
    \item \textbf{Trend analysis}: Mann-Kendall test for turnout decay
\end{itemize}

Significance level $\alpha = 0.05$ with Bonferroni correction.

\subsubsection{Effect Size}

Cohen's $d$ for pairwise comparisons:
\begin{equation}
d = \frac{\bar{x}_1 - \bar{x}_2}{s_{\text{pooled}}}
\end{equation}

\subsection{Validation}

\subsubsection{Internal Validity}
\begin{itemize}
    \item Deterministic seeding ensures reproducibility
    \item Multiple runs per configuration establish distributions
    \item Systematic parameter variation isolates effects
\end{itemize}

\subsubsection{External Validity}
\begin{itemize}
    \item Agent models are simplified; real voters have complex motivations
    \item Fatigue model parameters derived from intuition, not empirical calibration
    \item No adversarial behavior modeled in these experiments
\end{itemize}
